\section{Exponential cancellation at proportional harmonic depth}
\label{sec:saddle}

The algebraic results above organize special arguments. We now turn to an
analytic question left open by the earlier alternating-harmonic report
\cite{repoHarmonic}: how rapidly does its normalized sum tend to zero
below the first-term transition? The case of proportional outer exponent
and harmonic depth has an explicit rate function. This is a continuation
of that report's results, not another proof of its already established
reflection identity.

For integers $h,n\geq0$, define
\[
 e_h(n)=[u^h]\prod_{j=1}^n(1+u/j),\qquad
 T_{p,h}(a)=\sum_{n\geq0}\frac{(-1)^ne_h(n)}{(n+a)^p},
 \quad a,p>0.
\]
Thus $e_0(n)=1$, $e_1(n)=H_n$, and
$e_2(n)=(H_n^2-H_n^{(2)})/2$. Ordinary convergence holds for $p>0$;
absolute convergence holds for $p>1$. The Gaussian normalization is
$S_{p,h}=2^{-p}T_{p,h}(1/2)$. The positive Mellin integral
\begin{equation}\label{eq:positive-mellin}
 T_{p,h}(a)=\frac{(-1)^h}{h!\,\Gamma(p)}
 \int_0^\infty
 \frac{x^{p-1}e^{-ax}\log^h(1+e^{-x})}{1+e^{-x}}\,dx
\end{equation}
follows from
\[
 \sum_{n\geq0}e_h(n)z^n
 =\frac{(-\log(1-z))^h}{h!(1-z)}
\]
by inserting an Abel factor, using the gamma integral, and taking its
limit. The exponential decay at infinity and the bound
$\log(1+e^{-x})\leq\log2$ near zero justify the integral limit.
The convergence and Abel-limit argument are also proved in
\cite[Section 3]{repoHarmonic}. In particular,
$(-1)^hT_{p,h}(a)>0$.

Put
\begin{equation}\label{eq:R-definition}
 R_{p,h}(a)=(-1)^h h!(h+a)^p T_{p,h}(a).
\end{equation}
The earlier report proves $0<R_{p,h}(a)<1$ and identifies the threshold
$p\sim h\log h$. When $p/h-\log h\to-\infty$, it proves that
$R_{p,h}(a)\to0$. The following theorem quantifies this statement when
$p/h$ remains in a fixed positive compact interval.

\subsection{The unique saddle and the rate function}

Write
\begin{equation}\label{eq:L-w}
 L(x)=\log(1+e^{-x}),\qquad
 w(x)=\frac1{(1+e^x)L(x)}=-\frac{d}{dx}\log L(x),\quad x>0.
\end{equation}

\begin{lemma}\label{lem:w}
The function $w$ is strictly increasing, with
\[
 \frac1{2\log2}<w(x)<1.
\]
For every $\alpha>0$ there is exactly one $x_\alpha>0$ satisfying
\begin{equation}\label{eq:saddle-equation}
 x_\alpha w(x_\alpha)=\alpha.
\end{equation}
Moreover $\alpha<x_\alpha<2\alpha\log2$.
\end{lemma}
\begin{proof}
Set $t=e^{-x}\in(0,1)$. Then
$w=t/((1+t)\log(1+t))$ and
\[
 \frac{d}{dt}\log w
 =\frac{\log(1+t)-t}{t(1+t)\log(1+t)}<0.
\]
Consequently $w$ increases with $x$. Its endpoint limits give the bounds.
The function $xw(x)$ has derivative $w+xw'>0$, tends to zero at zero,
and tends to infinity at infinity. This proves existence, uniqueness,
and the stated enclosure.
\end{proof}

For $x=x_\alpha$, set
\begin{align}
 \Phi_\alpha(y)&=\alpha\log y+\log L(y),\label{eq:phase}\\
 A_\alpha&=\frac{\alpha}{x^2}+w'(x)>0,\label{eq:A}\\
 I(\alpha)&=-\alpha\log(x/\alpha)-\alpha-\log L(x),\label{eq:rate}\\
 C(\alpha,a)&=
 \frac{\sqrt{\alpha/A_\alpha}}{x(1+e^{-x})}
 \exp\bigl(a(\alpha-x)\bigr).\label{eq:C}
\end{align}
All of these functions are real analytic for $\alpha>0$.

\begin{theorem}[Proportional-depth asymptotics]\label{thm:saddle}
Let $a$ and $\alpha$ range over fixed compact subsets of $(0,\infty)$,
and put $p=\alpha h$. As the integer $h\to\infty$,
\begin{equation}\label{eq:saddle-leading}
 R_{\alpha h,h}(a)
 =C(\alpha,a)e^{-hI(\alpha)}\bigl(1+O(h^{-1})\bigr),
\end{equation}
uniformly on those compact sets. In fact
\begin{equation}\label{eq:saddle-second}
 R_{\alpha h,h}(a)
 =C(\alpha,a)e^{-hI(\alpha)}
 \left(1+\frac{D(\alpha,a)}h+O(h^{-2})\right),
\end{equation}
where, with $x=x_\alpha$, $A=A_\alpha$, and
$g(y)=e^{-ay}/(y(1+e^{-y}))$,
\begin{align}
 D(\alpha,a)={}&
 \frac{g''(x)}{2Ag(x)}
 +\frac{\Phi_\alpha'''(x)g'(x)}{2A^2g(x)}
 +\frac{\Phi_\alpha''''(x)}{8A^2}
 +\frac{5\Phi_\alpha'''(x)^2}{24A^3}
 -\frac{\alpha a^2}{2}-\frac1{12\alpha}.
 \label{eq:D}
\end{align}
Every fixed further order in $h^{-1}$ exists and is effectively obtained
from derivatives of $L$ and $g$ at $x_\alpha$ and the Stirling series.
The remainder at every fixed order is uniform on the same compact sets.
\end{theorem}

\begin{proof}
Equation \eqref{eq:positive-mellin} gives
\[
 R_{\alpha h,h}(a)
 =\frac{(h+a)^{\alpha h}}{\Gamma(\alpha h)}
 \int_0^\infty e^{h\Phi_\alpha(y)}g(y)\,dy.
\]
We have
$\Phi_\alpha'(y)=\alpha/y-w(y)$ and
$\Phi_\alpha''(y)=-\alpha/y^2-w'(y)<0$.
The sole critical point is therefore a strict global maximum.
Its location stays in a fixed compact subinterval of $(0,\infty)$ as
$\alpha$ varies in the prescribed compact set, and its curvature is
bounded away from zero.

Here are explicit tail controls for the application of Laplace's method
\cite[\S2.3(iii)]{DLMFLaplace}
\cite[Section 2.3(iii)]{DLMFLaplace}. Near zero,
$\exp(h\Phi_\alpha(y))\leq y^{h\alpha}(\log2)^h$,
while $g(y)=O(y^{-1})$ uniformly in $a$. Choosing the endpoint interval
small enough makes its integral exponentially smaller than the maximum,
uniformly in $\alpha$. At infinity, $L(y)<e^{-y}$ implies
$\Phi_\alpha(y)<\alpha\log y-y$; this supplies an integrable
exponential bound uniform in the same compact parameters. On the
remaining compact set outside a fixed neighborhood of $x_\alpha$,
strict concavity and compactness give a uniform negative gap from the
maximum. These observations justify discarding both tails with an
exponentially small relative error.

In the neighborhood, substitute
$y=x+v/\sqrt h$. Taylor expansion, followed by Gaussian integration,
gives
\begin{align*}
 \int_0^\infty e^{h\Phi_\alpha(y)}g(y)\,dy
 ={}&e^{h\Phi_\alpha(x)}g(x)
 \sqrt{\frac{2\pi}{hA}}\left(1+\frac Eh+O(h^{-2})\right),\\
 E={}&\frac{g''}{2Ag}
 +\frac{\Phi'''g'}{2A^2g}
 +\frac{\Phi''''}{8A^2}
 +\frac{5(\Phi''')^2}{24A^3},
\end{align*}
with all derivatives evaluated at $x$. Odd Gaussian moments vanish;
the even moments $1/A$, $3/A^2$, and $15/A^3$ produce precisely these
four terms. The analytic functions have uniformly bounded derivatives
on a common neighborhood of the saddle. Taylor remainders there,
together with the tail bounds just given, prove the claimed uniform
error and also permit expansion to any prescribed order.

Finally Stirling's formula, uniformly for $\alpha$ in its compact set,
and $\alpha h\log(1+a/h)=\alpha a-\alpha a^2/(2h)+O(h^{-2})$
give
\[
 \frac{(h+a)^{\alpha h}}{\Gamma(\alpha h)}
 =\sqrt{\frac{\alpha h}{2\pi}}
 e^{h\alpha(1-\log\alpha)+a\alpha}
 \left(1-\frac{\alpha a^2/2+1/(12\alpha)}h+O(h^{-2})\right).
\]
Multiplying the two expansions proves
\eqref{eq:saddle-leading}--\eqref{eq:D}.
\end{proof}

\begin{remark}[Integer exponents and rounding]
For integer $p$, set $\alpha=p/h$ in the theorem. If instead one fixes
$\alpha_0$ and rounds $\alpha_0h$ to $p$, the $O(1/h)$ change of
$\alpha$ must remain inside $e^{-hI(\alpha)}$: replacing it by
$e^{-hI(\alpha_0)}$ changes the leading constant by a generally
nontrivial bounded factor. This is a practical distinction in numerical
comparisons.
\end{remark}

\subsection{Monotonicity, convexity, and matching limits}

\begin{proposition}\label{prop:rate-properties}
The rate $I$ is positive, strictly decreasing, and strictly convex on
$(0,\infty)$. More explicitly,
\begin{equation}\label{eq:rate-derivatives}
 I'(\alpha)=\log\frac{\alpha}{x_\alpha}=\log w(x_\alpha)<0,
 \qquad
 I''(\alpha)=\frac{w'(x_\alpha)}
 {w(x_\alpha)(w(x_\alpha)+x_\alpha w'(x_\alpha))}>0.
\end{equation}
Its endpoint limits are
\[
 I(0+)=-\log\log2,\qquad I(+\infty)=0.
\]
As $\alpha\to\infty$,
\begin{align}
 x_\alpha&=\alpha+\frac{\alpha}{2}e^{-\alpha}
             +O(\alpha^2e^{-2\alpha}),\label{eq:x-large}\\
 I(\alpha)&=\frac12e^{-\alpha}
      -\frac{3\alpha+5}{24}e^{-2\alpha}
      +O(\alpha^2e^{-3\alpha}),\label{eq:I-large}\\
 C(\alpha,a)&=1-
 \left(\frac{(a+1/2)\alpha}{2}+1\right)e^{-\alpha}
 +O_a(\alpha^2e^{-2\alpha}).\label{eq:C-large}
\end{align}
Also $C(\alpha,a)\to1/2$ as $\alpha\downarrow0$.
\end{proposition}
\begin{proof}
Differentiating the optimized value in \eqref{eq:rate} cancels its
$x_\alpha'$ terms by \eqref{eq:saddle-equation}. This gives the first
identity in \eqref{eq:rate-derivatives}. Differentiating
$\alpha=xw(x)$ gives the second. Positivity also follows directly:
$\log L(y)<-y$, so
\[
 \max_{y>0}\{\alpha\log(y/\alpha)+\alpha+\log L(y)\}
 <\max_{y>0}\{\alpha\log(y/\alpha)+\alpha-y\}=0.
\]
The strict inequality is valid because the first maximum is attained;
at its maximizing point the second expression is at most zero.

The small-$\alpha$ limits follow from
$x_\alpha/\alpha\to2\log2$ and
$A_\alpha\sim\alpha/x_\alpha^2$. At infinity, put $t=e^{-y}$ and use
\[
 \log L(y)=-y-\frac t2+\frac{5t^2}{24}+O(t^3),
 \qquad
 w(y)=1-\frac t2+\frac{5t^2}{12}+O(t^3).
\]
Substitution in \eqref{eq:saddle-equation} first gives
\eqref{eq:x-large}. Taylor expansion of \eqref{eq:rate} around
$y=\alpha$, retaining the quadratic displacement, gives
\eqref{eq:I-large}. Expansion of \eqref{eq:C} and
$A_\alpha=\alpha/x_\alpha^2+w'(x_\alpha)$ gives
\eqref{eq:C-large}. The remainders stated above follow from analyticity
of $\log(\log(1+t)/t)$ for $t$ near zero, with
$x_\alpha-\alpha=O(\alpha e^{-\alpha})$.
\end{proof}

The rate formula explains the scale of cancellation even when the terms
of the original alternating sum are much larger than its value. For
example, $I(1)=0.15091647898554\ldots$ and
$C(1,1/2)=0.66325765726435\ldots$. Thus
\[
 R_{h,h}(1/2)\sim0.66325765726435\ldots\,
 e^{-0.15091647898554\ldots h}.
\]
This is a proved exponential rate, rather than just a limit of zero.

The large-$\alpha$ expansions also explain why the transition parameter
in the previous report was $\lambda=he^{-\alpha}/2$. Formally putting
$\alpha=\log h+c$ in $hI(\alpha)$ gives the leading value
$e^{-c}/2$. This agrees with the independently proved transition
$R\to\exp(-e^{-c}/2)$ in \cite{repoHarmonic}. Such matching is a
consistency calculation: the compact-$\alpha$ theorem proved here does
\emph{not} by itself establish a uniform error for growing $\alpha$.
Obtaining a single effective formula across that expanding range
remains a precise further problem.

\subsection{A probability law at the saddle}

There is a useful formulation independent of the alternating series.
Normalize the positive integrand in \eqref{eq:positive-mellin}, with
$p=\alpha h$, to a probability density and call the resulting random
variable $X_h$. Then
\begin{equation}\label{eq:saddle-clt}
 \sqrt{hA_\alpha}(X_h-x_\alpha)
 \ \Longrightarrow\ N(0,1).
\end{equation}
Indeed, the same change of variables used in the proof of
Theorem~\ref{thm:saddle}, its uniform Gaussian domination on a shrinking
neighborhood, and its exponentially negligible tails prove convergence
against every bounded continuous test function. The associated rate
for $X_h$ is
$\Phi_\alpha(x_\alpha)-\Phi_\alpha(y)$.
This identifies which part of the positive integral carries the very
small alternating sum and gives a stable center for numerical quadrature.

